\newglossaryentry{AsOp}{name={$\AsOp$}, description={``{\tt as}'' operator to bind a new variable to a structure}}

\chapter{Introduction}\label{ch:intro}
\minitoc

Pattern matching is a programming language mechanism to test a structure against a pattern at run-time.
It has two main purposes, matching and decomposition.
Matching is to determine whether the tested structure has the particular components organized in the particular way specified in the pattern.
Once the matching is successful, the individual components of the structure can be accessed through new variable bindings introduced in the pattern.
Pattern matching supersedes conditional statements in many occasions by stronger specificity, better integrity and higher abstraction,
where conditions are organized as a whole structure that is more intuitive and efficient, easier to comprehend, and less complicated.

\begin{figure}[t]
\centering
$$
\Br(\Br(x, 2, y \mathbin{\mgls{AsOp}} \Br(\Lf, 3, \Lf)), 4, z)
$$
\begin{tikzpicture}[default,->]
\node[nvertex] {$\Br$} [level distance=30pt,sibling distance=40pt]
  child {node[nvertex] (w) {$\Br$}
    child {node[vertex] (x) {}}
    child {node[nvertex] {$2$}}
    child {node[nvertex] (y) {$\Br$}
      child {node[nvertex] {$\Lf$}}
      child {node[nvertex] {$3$}}
      child {node[nvertex] {$\Lf$}}}}
  child {node[nvertex] {$4$}}
  child {node[vertex] (z) {}};

\draw[var] (x.west) -- node[pos=1.0,left] {$x$} +(left:20pt);
\draw[var] (y.east) -- node[pos=1.0,right] {$y$} +(right:20pt);
\draw[var] (z.east) -- node[pos=1.0,right] {$z$} +(right:20pt);
\end{tikzpicture}
\caption{Syntax tree of an algebraic pattern}\label{fig:al-pattern}
\end{figure}


\section{Algebraic Pattern Matching}

Pattern matching is already common in functional programming.
Patterns in the functional world follow the scheme of algebraic data types, or variants, where data objects are generated by hierarchical constructor applications.
The patterns are therefore \emph{partial} applications, structured as abstract syntax trees of function invocations.
\autoref{fig:al-pattern} shows the syntax tree of a typical binary search tree pattern in a functional language,
where $x$, $y$ and $z$ provide new bindings of variables to the nodes in the application tree.
Such a pattern follows the hierarchically structured algebraic constructors, that is simple and easy to match with very clear semantics.
That's why pattern matching is widely used in the very early ages of functional languages.

When pattern matching is migrated to \gls{OO} languages, subtyping between classes must be brought into consideration.
A great benefit gained is the ability to perform type tests with variable bindings in an atomic way,
that is much safer and eliminates the need of typecasts.
However, the algebraic nature of patterns does not change very much.
An additional layer is commonly introduced to bridge the difference between the algebraic view~\cite{burton1993pattern}
and the underlying pointer structure of an object, for example, the \emph{case class} in Scala~\cite{odersky2021scala}.
Although the case class constructors create objects with stored structures identical to normal objects,
those stored structures are in fact irrelevant.
The Scala \emph{extractors} can completely virtualize the case class constructors and abstract the stored structures away.
Such pattern matching matches an object against it abstract meaning expressed by the algebraic view rather than the underlying stored structure.
Therefore, the semantics of patterns is still in the functional world, with the benefit of higher level abstraction,
and the limitation of unable to examine the stored relations of data objects.

\begin{figure}[t]
\centering
\begin{tikzpicture}[default,->]
\node[root] {} [level distance=30pt,sibling distance=60pt]
  child {node[vertex] (w) {}
    child {node[vertex] (x) {} edge from parent node[below,pos=0.4] {$l$}}
    child {node[vertex] (y) {} edge from parent[draw=none]}
    edge from parent node[above] {$l$}}
  child {node[vertex] (z) {} edge from parent node[above] {$r$}};

\draw (w) to[bend left] node[above,pos=0.6] {$r$} (y);
\draw (y) to[bend left] node[below] {$l$} (w);
\draw (y) -- node[above,pos=0.4] {$r$} (z);

\draw[var] (w) -- node[pos=1.0,left] {$w$} +(135:20pt);
\draw[var] (x) -- node[pos=1.0,left] {$x$} +(left:20pt);
\draw[var] (y) -- node[pos=1.0,right] {$y$} +(-45:20pt);
\draw[var] (z) -- node[pos=1.0,right] {$z$} +(right:20pt);
\end{tikzpicture}
\caption{Object with aliased and circular relations}\label{fig:obj-pattern}
\end{figure}

The decision to put {\OO} pattern matching in a functional way has its reasons.
Besides that the semantics of algebraic patterns is simpler and more familiar to programmers who already know functional pattern matching,
the underlying pointer structures of objects can have aliases and cycles, that are more complicated than trees.
\autoref{fig:obj-pattern} shows an object with aliased and circular relations.
Furthermore, when considering the conjunction and disjunction of multiple patterns with subtyping,
a type that is compatible with many types must be resolved.
Nevertheless, there's a need to expose and explore the concrete stored structures of objects when writing internal implementations for {\OO} methods.
This is often the case where pattern matching is used, since diverging processing of variants externally in {\OO} paradigm usually goes through the subtyping polymorphism.
Also, expressing the relations between objects explicitly in the stored structures is a key characteristic of {\OO} programs.
Therefore, by addressing these challenges, we propose a pattern matching extension to general {\OO} languages
that directly operates on the pointer structures of objects commonly found in most language implementations.

\section{Graph-based Pattern Matching}

We base the theoretical foundation of our extension on the well established object graph model~\cite{ke2009graph}.
An object in this model is a rooted and directed graph with labeled edges.
The root represents the object with outgoing edges pointing to components.
The edges are labeled with the attributes of the source node.
Local variables mark the edges from the stack to roots of objects.
\autoref{fig:junction} shows an object graph on the upper-left, where local variable $o$ leads to the object, which has attributes $o.a$, $o.b$, $o.c$, $o.d$ and $o.e$.
Among the attributes, $o.a$ and $o.b$ alias each other and point to the same object.
We omit the stack in the shown graphs, since it's somewhat irrelevant to pattern matching. Those edges without a source node are coming from the stack.
In our pattern matching graph model, both objects and patterns are represented by the object graphs, which are collectively called pattern graphs in the context.
The major difference is that a pattern is not a complete object, some or all of the attributes can be omitted.

\begin{figure}[t]
\centering
\begin{tikzpicture}[default,->,x=35pt,y=35pt]
%% o

\node[root] (r) at(0,0) {};
\node[vertex,label=below:{$w_1$}] (b) at ($(r)+(-1,-1)$) {};
\node[vertex,label=below:{$w_2$}] (c) at ($(r)+(0,-1)$) {};
\node[vertex,label=below:{$w_3$}] (d) at ($(r)+(1,-1)$) {};
\node[vertex] (e) at ($(r)+(1,0)$) {};

\draw (r) to[bend right=45] node[left,pos=0.75] {$a$} (b);
\draw (r) -- node[below] {$b$} (b);
\draw (r) -- node[right] {$c$} (c);
\draw (r) -- node[above,pos=0.75] {$d$} (d);
\draw (r) -- node[above] {$e$} (e);

\draw[var] (r) -- node[pos=1.0,left] {$o$} +(135:0.67);

\draw[-,densely dashed] (-1.25,-1.825) -- +(8,0);
\draw[-,densely dashed] (4.125,0.25) -- +(0,-4.5);

%% p1

\node[root] (r) at (0,-2.75) {};
\node[vertex,label=below:{$u_1$}] (a) at ($(r)+(-1,0)$) {};
\node[vertex,label=below:{$u_2$}] (c) at ($(r)+(0,-1)$) {};
\node[vertex] (d) at ($(r)+(1,-1)$) {};

\draw (r) -- node[above] {$a$} (a);
\draw (r) -- node[left] {$c$} (c);
\draw (r) -- node[above,pos=0.75] {$d$} (d);

\draw[var] (r) -- node[pos=1.0,right] {$p_1$} +(45:0.67);
\draw[var] (c) -- node[pos=1.0,left] {$x$} +(left:0.67);

%% p2

\node[root] (r) at (2.75,-2.75) {};
\node[vertex,label=below:{$u_3$}] (b) at ($(r)+(-1,-1)$) {};
\node[vertex,label=below:{$u_4$}] (c) at ($(r)+(0,-1)$) {};
\node[vertex] (e) at ($(r)+(1,0)$) {};

\draw (r) -- node[above,pos=0.75] {$b$} (b);
\draw (r) -- node[left] {$c$} (c);
\draw (r) -- node[above] {$e$} (e);

\draw[var] (r) -- node[pos=1.0,left] {$p_2$} +(135:0.67);

%% p3

\node[root] (r) at (5.5,-2.75) {};
\node[vertex,label=below:{$u_5$}] (b) at ($(r)+(-1,-1)$) {};
\node[vertex,label=below:{$u_6$}] (c) at ($(r)+(0,-1)$) {};
\node[vertex,label=below:{$u_7$}] (d) at ($(r)+(1,-1)$) {};

\draw (r) -- node[above,pos=0.75] {$b$} (b);
\draw (r) -- node[left] {$c$} (c);
\draw (r) -- node[above,pos=0.75] {$d$} (d);

\draw[var] (r) -- node[pos=1.0,right] {$p_3$} +(45:0.67);

%% p4

\node[root] (r) at (5.5,0) {};
\node[vertex] (b) at ($(r)+(-1,-1)$) {};
\node[vertex,label=below:{$u_8$}] (d) at ($(r)+(1,-1)$) {};

\draw (r) -- node[above,pos=0.75] {$b$} (b);
\draw (r) to[bend right] node[below] {$c$} (d);
\draw (r) to[bend left] node[above] {$d$} (d);

\draw[var] (r) -- node[pos=1.0,left] {$p_4$} +(135:0.67);
\draw[var] (d) -- node[pos=1.0,right] {$y$} +(45:0.67);
\end{tikzpicture}
\caption{Junction of multiple patterns}\label{fig:junction}
\end{figure}

Given the notion of pattern graphs, our pattern matching is intuitive.
A pattern graph matches another when the one \emph{covers} the other.
To define the covering between two graphs, we need to figure out the correspondence between the nodes of the two graphs.
Beginning with the roots, the nodes led to by the edges with the same label correspond to each other.
Cumulatively, the label sequence from the root to a node is the incoming path, or the \emph{trace}, of a node.
Although a single path is sufficient to find a node, we must take into account the whole set of incoming paths to a node to obtain the graph structure.
The idea to identify a node by how it can be reached is called the trace model~\cite{hoare1999trace},
where an isomorphism can be built between a node and the set of its incoming paths.
The trace model is crucially concise in reasoning about equivalent nodes in different graphs, and plays extremely well with aliases and cycles.
There are two aspects in determination of the covering --- the layout test and the type test.
First, a node $w_1$ covers another node $u_1$ when the incoming paths of $w_1$ include the incoming paths of $u_1$.
To separate the case of aliased nodes, a node can only cover one node in the other graph.
As shown in \autoref{fig:junction}, $w_1$ covers $u_1$ for $\sett{[a],[b]} \supset \sett{[a]}$,
on the contrary, $u_8$ covers either $u_6$ or $u_7$, but not both.
Second, the type $\tau(w_1)$ must be a subtype of $\tau(u_1)$.
In {\OO} languages, subtyping is a preorder, i.e., reflexive and transitive~\cite{abadi1996theory,liskov1994behavioral}.
We only stick with this abstraction in order to make the pattern matching general.
When all the nodes of a pattern graph are covered by the nodes of the other, we say this pattern is matched by the other.
In terms of subtyping, the match relation is covariant.

It's getting more interesting when we look further into the match against multiple patterns.
One object can be tested to match all of the patterns, the conjunction, or at least one of the patterns, the disjunction.
For example, in \autoref{fig:junction}, object $o$ matches all of $p_1$, $p_2$ and $p_3$, and one of $p_3$ and $p_4$.
While the tests of multiple patterns are straightforward,
the bindings of variable $x$ in the conjunction, and $y$ in the disjunction are more subtle.
For conjunction, if there exists a node $w_2$ that can cover $u_2$, $u_4$ and $u_6$, $\tau(w_2)$ must be a subtype of all $\tau(u_2)$, $\tau(u_4)$ and $\tau(u_6)$,
therefore it's possible to assign a better type to variable $x$ than just $\tau(u_2)$.
Also, there exist invalid conjunctions which no graph can possibly match.
In \autoref{fig:junction}, $p_3$ and $p_4$ cannot be matched by one graph. We need to detect and reject invalid conjunctions.
For disjunction, we don't have conflicts between the multiple patterns.
However, without knowing which pattern will be eventually matched at run-time, variable bindings are more complicated.
For example, what is the type of variable $y$ in the disjunction of $p_3$ and $p_4$, $\tau(u_8)$, $\tau(u_6)$ or $\tau(u_7)$?
If object $o$ matches the disjunction, for it matches $p_3$, which node do we bind to $y$, $w_2$ or $w_3$?

\section{Contribution and Outline}
With the brief discussion of the problems and ideas, below lists the challenges we are going to address in this work.
\begin{enumerate}
\item
We build a formal base for the pattern graphs with subtyping.
This includes the definitions, morphisms, operations, match relations and algorithms for matching.
In addition, we design a language to construct the pattern graphs, in particular with aliases and cycles,
as an extension to existing {\OO} languages.

\item
We study the compatibility of the pattern matching with the current type-checking and run-time environment,
and find a way for smooth integration.
We present the type-checking rules, operational semantics and reason about the type safety.

\item
We explore the conjunction and disjunction of multiple patterns,
and give our resolution, design decision and algorithms to handle the relatively complex cases at compile-time.
\end{enumerate}

\subsection{Current Publication}\label{sec:publication}

\begin{BlindRegion}

\begin{enumerate}[label=(C\arabic*)]
  \item \bibentry{ke2009graph}
\end{enumerate}

\begin{enumerate}[label=(J\arabic*)]
  \item \bibentry{ke2012rcos}
  \item \bibentry{ke2013graph}
  \item \bibentry{ke2021pattern}
\end{enumerate}

\end{BlindRegion}

\subsection{Thesis Outline}

The rest of the thesis is outlined as follows.
\autoref{ch:related} summarizes the work in object graphs and pattern matching, from the perspective of {\OO} programming.
\autoref{ch:graph} defines the formal base of pattern matching on object graphs, including pattern conjunction and disjunction,
with justifications and constructive algorithms.
\autoref{ch:language} presents the language extension for the graph-based pattern matching in abstract syntax,
and provides the type-checking rules and operational semantics, together with the proof of type safety.
\autoref{ch:discussion} discusses the restriction and limitation of the theory, and the potential solutions.
Finally, \autoref{ch:conclusion} concludes the work.

\endinput